\subsubsection{Frontend}

\begin{figure}[H]
	\centering 
	\includegraphics[width=1\textwidth]{SpecificaTecnica/Autenticazione/UMLAutenticazioneFrontend}
	\caption{Diagramma delle classi per Autenticazione e Registrazione: Frontend}
\end{figure}

\subsubsubsection{FormView}

\begin{figure}[H]
	\centering 
	\includegraphics[width=0.4\textwidth]{SpecificaTecnica/Autenticazione/FormView}
	\caption{Autenticazione e Registrazione: FormView}
\end{figure}

\paragraph{Descrizione}
\texttt{FormView} è il componente React$^G$ che costituisce il livello di \textit{View} nel pattern MVVM.\\
Si occupa esclusivamente della presentazione grafica del form di autenticazione o registrazione, delegando tutta la logica al \texttt{FormViewModel} ricevuto tramite \texttt{FormProps}.

\paragraph{Attributi}
\begin{itemize}
	\item \texttt{visibleFields: Record} - attributo privato di tipo \texttt{Record<string, boolean>}.
	Gestisce la visibilità dei campi password: per ciascun campo di tipo \texttt{password}, traccia se il testo è attualmente visibile o nascosto.
	Viene aggiornato dalla funzione \texttt{toggleVisibility} al click sull'icona dell'occhio.
\end{itemize}

\paragraph{Metodi e funzioni}
\begin{itemize}
	\item \texttt{Form(f: FormProps)} - funzione componente React$^G$ (funzione pubblica).
	Riceve le props di tipo \texttt{FormProps} e restituisce il markup JSX del form.
	Itera sulla lista \texttt{fields} per renderizzare dinamicamente i campi di input, utilizza \texttt{vm.values} e \texttt{vm.fieldErrors} per il binding dei dati, e collega gli handler \texttt{vm.handleChange}, \texttt{vm.handleBlur} e \texttt{vm.handleSubmit} agli eventi del form. Per i campi di tipo password gestisce la visibilità tramite \texttt{visibleFields} e delega la visualizzazione degli errori globali al componente \texttt{FormErrors}.
\end{itemize}


\subsubsubsection{FormViewModel}

\begin{figure}[H]
	\centering 
	\includegraphics[width=0.6\textwidth]{SpecificaTecnica/Autenticazione/FormViewModel}
	\caption{Autenticazione e Registrazione: FormViewModel}
\end{figure}

\paragraph{Descrizione}
\texttt{FormViewModel} è un'interfaccia TypeScript che definisce il contratto del ViewModel nel pattern MVVM.\\
Espone lo stato del form e gli handler degli eventi che la View deve utilizzare.\\
La sua implementazione concreta è prodotta dall'hook \texttt{useFormViewModel}, che gestisce lo stato React$^G$ e la comunicazione con il Model.

\paragraph{Attributi}
\begin{itemize}
	\item \texttt{values: Record} - attributo pubblico di tipo \texttt{Record<string, string>}.
	Mappa chiave-valore contenente i valori correnti di ciascun campo del form, viene aggiornata ad ogni modifica$^G$ dell'utente.

	\item \texttt{fieldErrors: Record} - attributo pubblico di tipo \texttt{Record<string, string>}.
	Mappa chiave-valore contenente i messaggi di errore di validazione$^G$ per ciascun campo, popolata durante la validazione$^G$ on-blur e al momento del submit.

	\item \texttt{errors: string*} - attributo pubblico di tipo array di stringhe.
	Contiene i messaggi di errore globali restituiti dal backend$^G$ o generati durante l'invio del form.

	\item \texttt{loading: bool} - attributo pubblico booleano.
	Indica se è in corso una richiesta asincrona verso il backend.
	Viene utilizzato dalla View per disabilitare il pulsante di submit e mostrare un indicatore di caricamento.
\end{itemize}

\paragraph{Metodi e funzioni}
\begin{itemize}
	\item \texttt{handleChange(key: string, value: string): void} - metodo pubblico che aggiorna il valore del campo identificato da \texttt{key} all'interno di \texttt{values}.
	Viene invocato dalla View ad ogni evento \texttt{onChange} degli input.

	\item \texttt{handleBlur(key: string): void} - metodo pubblico che innesca la validazione$^G$ del singolo campo identificato da \texttt{key} nel momento in cui l'utente abbandona il campo.
	Aggiorna \texttt{fieldErrors} con l'eventuale errore riscontrato.

	\item \texttt{handleSubmit(e: React.SyntheticEvent): void} - metodo pubblico che gestisce la sottomissione del form intercettando l'evento del browser. Innesca la validazione$^G$ completa di tutti i campi e, se non vi$^G$ sono errori, invoca il metodo \texttt{submit} del Model delegando
	la chiamata API$^G$ al layer sottostante.
\end{itemize}

\subsubsubsection{FormModel}

\begin{figure}[H]
	\centering 
	\includegraphics[width=0.6\textwidth]{SpecificaTecnica/Autenticazione/FormModel}
	\caption{Autenticazione e Registrazione: FormModel}
\end{figure}

\paragraph{Descrizione}
\texttt{FormModel} è una classe astratta TypeScript che costituisce il livello di \textit{Model} nel pattern MVVM.\\
Incapsula la logica di validazione$^G$ dei campi del form e la logica di invio dei dati al backend$^G$ tramite il modulo \texttt{AuthAPI}.\\
Le classi concrete \texttt{LoginModel} e \texttt{RegisterModel} la estendono specializzandone il comportamento.\\
Solo \texttt{fields} e \texttt{submit} sono astratti: \texttt{validate} e \texttt{validateField} hanno invece un'implementazione concreta nella classe base.

\paragraph{Attributi}
\begin{itemize}
	\item \texttt{fields: FieldConfig*} - attributo pubblico astratto di tipo array di \texttt{FieldConfig}.
	Contiene la configurazione dei campi del form. Viene definito dalle sottoclassi concrete con i campi specifici per login o registrazione.

	\item \texttt{USERNAME\_REGEX: RegExp} - costante privata protetta di tipo \texttt{RegExp}.
	Espressione regolare utilizzata per la validazione$^G$ del formato del nome utente (tra 4 e 24 caratteri alfanumerici o underscore).

	\item \texttt{PASSWORD\_REGEX: RegExp} - costante privata protetta di tipo \texttt{RegExp}.
	Espressione regolare utilizzata per la validazione$^G$ della complessità della password (tra 8 e 24 caratteri, con almeno una maiuscola, una minuscola, un numero e un carattere speciale).
\end{itemize}

\paragraph{Metodi e funzioni}
\begin{itemize}
	\item \texttt{validate(values: Record<string, string>): Record<string, string>}\\
	Metodo pubblico concreto, riceve una mappa dei valori correnti del form e restituisce una mappa di errori.\\
	Le sottoclassi possono sovrascriverlo per estendere le regole di validazione$^G$ specifiche.

	\item \texttt{validateField(key: string, value: string): string}\\
	Metodo pubblico concreto, valida il singolo campo identificato da \texttt{key} utilizzando \texttt{validate} sull’intero set di valori aggiornato con \texttt{value}.\\
	Restituisce il messaggio di errore per il campo richiesto, oppure una stringa vuota se il campo è valido.

	\item \texttt{submit(values: Record<string, string>): Promise<SubmitResult>}\\
	Metodo pubblico astratto asincrono, costruisce il DTO appropriato a partire dai valori del form e delega la chiamata HTTP al modulo \texttt{AuthAPI}.\\ 
	Restituisce una \texttt{Promise<SubmitResult>} che il ViewModel utilizza per aggiornare lo stato della View in base all’esito.

	\item \texttt{login(dto: LoginDto): Promise<AuthResponse>}\\
	Metodo privato protetto, invia una richiesta di login tramite \texttt{authFetch}.

	\item \texttt{register(dto: RegisterDto): Promise<AuthResponse>}\\
	Metodo privato protetto, invia una richiesta di registrazione tramite \texttt{authFetch}.

	\item \texttt{authFetch(endpoint: string, body: object): Promise<AuthResponse>}\\
	Metodo privato protetto, incapsula la logica comune di chiamata all’API di autenticazione, gestendo eventuali errori di rete o risposte non valide.
\end{itemize}


\subsubsubsection{RegisterModel}

\begin{figure}[H]
	\centering 
	\includegraphics[width=0.4\textwidth]{SpecificaTecnica/Autenticazione/RegisterModel}
	\caption{Autenticazione e Registrazione: RegisterModel}
\end{figure}

\paragraph{Descrizione}
\texttt{RegisterModel} è una classe concreta che estende \texttt{FormModel}, specializzandone il comportamento per il caso d'uso$^G$ della registrazione di un nuovo utente.\\
Inizializza i campi del form con la configurazione specifica per la registrazione (username, email, password, conferma password) e sovrascrive \texttt{validate} per aggiungere la validazione$^G$ del formato email e la verifica$^G$ della corrispondenza delle due password.

\paragraph{Attributi}
\begin{itemize}
	\item \texttt{EMAIL\_REGEX: RegExp} - attributo privato di tipo \texttt{RegExp}. 
	Espressione regolare utilizzata per la validazione$^G$ del formato dell'indirizzo email inserito dall'utente durante la registrazione.
\end{itemize}

\paragraph{Metodi e funzioni}
\texttt{RegisterModel} sovrascrive il metodo \texttt{validate} ereditato da \texttt{FormModel}, estendendolo con la validazione$^G$ del formato username tramite \texttt{USERNAME\_REGEX}, del formato email tramite \texttt{EMAIL\_REGEX}, della complessità della password tramite \texttt{PASSWORD\_REGEX} e della corrispondenza tra i campi \texttt{password} e \texttt{confirmPwd}.


\subsubsubsection{LoginModel}

\begin{figure}[H]
	\centering 
	\includegraphics[width=0.4\textwidth]{SpecificaTecnica/Autenticazione/LoginModel}
	\caption{Autenticazione e Registrazione: LoginModel}
\end{figure}

\paragraph{Descrizione}
\texttt{LoginModel} è una classe concreta che estende \texttt{FormModel}, specializzandone il comportamento per il caso d'uso$^G$ dell'autenticazione di un utente esistente.\\
Inizializza i campi del form con la configurazione specifica per il login (username e password) e utilizza le regole di validazione$^G$ ereditate da \texttt{FormModel} senza aggiungerne di nuove.

\paragraph{Attributi}
\texttt{LoginModel} non definisce attributi aggiuntivi rispetto alla classe padre \texttt{FormModel}.

\paragraph{Metodi e funzioni}
\texttt{LoginModel} eredita senza modifiche i metodi \texttt{validate} e \texttt{validateField} da \texttt{FormModel}. Il metodo \texttt{submit} viene specializzato per costruire un oggetto \texttt{LoginDto} con i campi username e password.
